% Séquence 3 : Fractions décimales et nombres décimaux
\setseqtitle{Fractions décimales et nombres décimaux}
\chapter{Fractions décimales et nombres décimaux}

\begin{objectifsbox}
	\textbf{Objectifs d'apprentissage de la séquence}
	\begin{itemize}
		\item Comprendre ce qu'est une fraction décimale et un nombre décimal
		\item Connaître les liens entre les unités de numération : unité, dizaine, centaine, millier, dixième, centième, millième
		\item Associer et utiliser différentes écritures d'un nombre décimal : écriture à virgule, fraction décimale, décomposition
		\item Comparer deux nombres décimaux
		\item Ordonner une liste de nombres décimaux
		\item Encadrer un nombre décimal par deux nombres entiers ou décimaux
		\item Intercaler un nombre décimal entre deux nombres décimaux
		\item Placer et repérer un nombre décimal sur une demi-droite graduée
	\end{itemize}
\end{objectifsbox}

\section{Des fractions décimales aux nombres décimaux}

Dans notre vie de tous les jours, nous rencontrons des nombres qui ne sont pas entiers.
Un paquet de gâteaux pèse 0,375 kg, un athlète court le 100 mètres en 9,85 secondes, et un article coûte 12,99 €.
Ces nombres à virgule sont appelés des \textbf{nombres \solution[3cm]{décimaux}}. Dans ce chapitre, nous allons découvrir comment ils fonctionnent, comment les écrire de différentes manières et comment les utiliser pour comparer, ranger et mesurer.
\subsection{Rappel : fractions simples}

\begin{rappelbox}
	\textbf{Fraction}
	
	Une \textbf{fraction} est un nombre qui s'écrit sous la forme $\dfrac{a}{b}$ où $a$ et $b$ sont des nombres entiers avec $b \neq \solution[0.5cm]{0}$.
	\begin{itemize}[label = \textbullet]
		\item $a$ est appelé le \textbf{\solution{numérateur}}
		\item $b$ est appelé le \textbf{\solution{dénominateur}}
	\end{itemize}
	
	La fraction $\dfrac{a}{b}$ représente le \solution[4cm]{quotient} de la division de $a$ par $b$.
\end{rappelbox}

\subsection{Fractions décimales}

\begin{definitionbox}
	\textbf{Fraction décimale}
	
	Une \textbf{fraction décimale} est une fraction dont le dénominateur est \solution{10}, \solution{100}, \solution{1~000}, etc.
	
	\textbf{Astuce :} Le dénominateur est un \solution{1} suivi uniquement de \solution{zéros}.
\end{definitionbox}

\begin{examplebox}
	\textbf{Exemples de fractions décimales :}
	$\dfrac{7}{10}$ (7 \solution{dixièmes}) ; $\dfrac{23}{100}$ (23 \solution{centièmes}) ; $\dfrac{145}{1~000}$ (145 \solution{millièmes}).
	\textbf{Contre-exemples (ce ne sont pas des fractions décimales) :}
	$\dfrac{2}{3}$ ; $\dfrac{5}{7}$ (leurs dénominateurs ne sont ni 10, ni 100, etc.).
\end{examplebox}

\subsection{Nombres décimaux}

\begin{definitionbox}
	\textbf{Nombre décimal}
	
	Un \textbf{nombre décimal} est un nombre qui peut s'écrire sous forme de \solution{fraction décimale}.
	On peut aussi l'écrire avec une \solution{virgule} : c'est l'écriture décimale.
	\textbf{Vocabulaire :}
	\begin{itemize}[label = \textbullet]
		\item La \textbf{partie \solution{entière}} est la partie avant la virgule.
		\item La \textbf{partie \solution{décimale}} est la partie après la virgule.
	\end{itemize}
\end{definitionbox}

\begin{examplebox}
	$0{,}7 = \solution{$\dfrac{7}{10}$}$ \quad ; \quad $12{,}56 = \solution{$\dfrac{1~256}{100}$}$ \quad ; \quad $3 = \solution{$\dfrac{30}{10}$}$ (les entiers sont des nombres décimaux !).\\
	La partie entière de 12,56 est \solution{12} et sa partie décimale est \solution{$0,56 = \dfrac{56}{100}}
\end{examplebox}

\section{Unités de numération}

\subsection{Tableau de numération}

\begin{center}
	\renewcommand{\arraystretch}{1.5}
	\begin{tabular}{|c|c|c|c|c|c|c|}
		\hline
		\multicolumn{3}{|c|}{\textbf{Partie entière}} & \multicolumn{4}{c|}{\textbf{Partie décimale}} \\
		\hline
		\textbf{Centaines} & \textbf{Dizaines} & \textbf{Unités} & \textbf{Dixièmes} & \textbf{Centièmes} & \textbf{Millièmes} & \textbf{Dix-millièmes} \\
		\hline
		$100$ & $10$ & $1$ & $\dfrac{1}{10}$ & $\dfrac{1}{100}$ & $\dfrac{1}{1~000}$ & $\dfrac{1}{10~000}$ \\
		\hline
		& 1 & 5 & 9 & 3 & 1 & \\
		\hline
	\end{tabular}
\end{center}
Le nombre représenté est \textbf{\solution{15,931}}.
Il se lit « \solution{quinze} unités et \solution{neuf cent trente-et-un} millièmes ».

\subsection{Décompositions d'un nombre décimal}

Un même nombre décimal peut être écrit de plusieurs façons. Il est important de savoir passer de l'une à l'autre pour mieux le comprendre.

Prenons l'exemple de \textbf{234,56}~:

\begin{center}
	\renewcommand{\arraystretch}{1.3}
\begin{tabularx}{\textwidth}{|>{\raggedright\arraybackslash}p{0.28\textwidth}|>{\raggedright\arraybackslash}p{0.32\textwidth}|>{\raggedright\arraybackslash}X|}
		\hline
		\textbf{Façon d'écrire} & \textbf{Exemple} & \textbf{Explication}\\
		\hline
		Écriture décimale & $234{,}56$ & C'est la forme habituelle.\\
		\hline
		Fraction décimale & \solution[1.5cm]{$\dfrac{23\,456}{100}$} & On enlève la virgule et divise par $100$ (2 chiffres après la virgule).\\
		\hline
		Partie entière + fraction décimale inférieur à 1 & \solution{$234 + \dfrac{56}{100}$} & 234 unités et 56 centièmes.\\
		\hline
		Décomposition additive par rang & \solution{$200 + 30 + 4 + 0{,}5 + 0{,}06$} & Chaque chiffre selon sa position.\\
		\hline
		Décomposition additive avec fractions & \solution{$200 + 30 + 4 + \dfrac{5}{10} + \dfrac{6}{100}$} & Chaque chiffre multiplié par sa fraction.\\
		\hline
		Décomposition canonique (ou selon les chiffres) & \solution{$(2\times 100)+(3\times 10)+(4\times 1)+(5\times 0{,}1)+(6\times 0{,}01)$} & Chiffre $\times$ valeur du rang.\\
		\hline
	\end{tabularx}
\end{center}

\vspace{1em}
Pour s'entraîner, choisis un autre nombre décimal et essaie de le décomposer de plusieurs façons !


\section{Les zéros dans la partie décimale}

\begin{proprietebox}
	\textbf{Propriété des zéros "inutiles"}
	
	Un nombre décimal ne change pas de valeur si on \solution{ajoute} ou si on \solution{supprime} des zéros à la \textbf{fin} de sa partie décimale.
	
	Cette propriété est très importante pour \solution{comparer} des nombres.
\end{proprietebox}

\begin{examplebox}
	\begin{itemize}[label = \textbullet]
		\item $3,5 = 3,5\solution{0} = 3,5\solution{00}$
		\item $12,900 = 12,9\solution{0} = 12,\solution{9}$
		\item Pour comparer $5,7$ et $5,67$, on peut écrire $5,7$ comme $5,7\solution{0}$. Il devient alors plus facile de voir que \solution{$5,70 > 5,67$}.
	\end{itemize}
\end{examplebox}

\section{Comparer et ranger des nombres décimaux}

\subsection{Définition et symboles}
Comparer deux nombres, c'est dire lequel est le plus \solution{grand}, le plus \solution{petit}, ou s'ils sont \solution{égaux}.
On utilise les symboles : \solution{$<$} (plus petit que), \solution{$>$} (plus grand que), \solution{$=$} (égal à).
\textbf{Astuce :} Le symbole pointe toujours vers le plus \solution{petit} nombre (la bouche du crocodile mange le plus grand !).

\subsection{Méthode de comparaison}

\begin{methodebox}
	\textbf{Pour comparer deux nombres décimaux :}
	\begin{enumerate}
		\item \textbf{On compare leurs parties \solution[4cm]{entières}.} Le nombre avec la plus grande partie entière est le plus grand.
		\item \textbf{Si les parties entières sont égales}, on compare leurs parties \solution[4cm]{décimales} chiffre par chiffre, de \solution[3cm]{gauche} à \solution[3cm]{droite} (dixièmes, puis centièmes...). Dès qu'un chiffre est plus grand que l'autre, on a trouvé le plus grand nombre.
	\end{enumerate}
\end{methodebox}

\begin{examplebox}
	\begin{itemize}
		\item Comparer 14,12 et 11,865 : La partie entière 14 est plus grande que 11, donc \solution[5cm]{$14,12 > 11,865$}.\\
		\item Comparer 27,28 et 27,6 : Les parties entières sont égales (27). On compare les dixièmes : 2 est plus petit que 6, donc \solution[5cm]{$27,28 < 27,6$}.
	\end{itemize}
\end{examplebox}

\clearpage
\subsection{Ranger une liste de nombres}
Ranger des nombres dans l'\textbf{ordre croissant}, c'est les classer du plus \solution{petit} au plus \solution{grand}.\\
Ranger des nombres dans l'\textbf{ordre décroissant}, c'est les classer du plus \solution{grand} au plus \solution{petit}.\\

\begin{examplebox}
	Ranger dans l'ordre croissant : 3,25 ; 2,36 ; 3,205 ; 3,3 ; 2,29.\\
	\textbf{Résultat :} \fcolorbox{black}{yellow!20}{\solution[10cm]{$2,29 < 2,36 < 3,205 < 3,25 < 3,3$}}
\end{examplebox}

\section{Encadrer et intercaler}

\subsection{Encadrer un nombre}
Encadrer un nombre, c'est le placer entre un nombre plus \solution{petit} et un nombre plus \solution{grand}.
On cherche souvent l'encadrement le plus précis possible.\\

\begin{methodebox}
	\begin{itemize}
		\item \textbf{Encadrer à l'unité :} entre deux \solution[4cm]{entiers} qui se suivent.
		\item \textbf{Encadrer au dixième :} entre deux nombres à \solution{un} chiffre après la virgule qui se suivent.
		\item \textbf{Encadrer au centième :} entre deux nombres à \solution{deux} chiffres après la virgule qui se suivent.
	\end{itemize}
\end{methodebox}

\begin{examplebox}
	Pour le nombre \textbf{7,384} :
	\begin{itemize}
		\item Encadrement à l'unité : \fcolorbox{black}{yellow!20}{\solution[5cm]{$7 < 7,384 < 8$}}\\
		\item Encadrement au dixième : \fcolorbox{black}{yellow!20}{\solution[5cm]{$7,3 < 7,384 < 7,4$}}\\
		\item Encadrement au centième : \fcolorbox{black}{yellow!20}{\solution[5cm]{$7,38 < 7,384 < 7,39$}}
	\end{itemize}
\end{examplebox}

\clearpage
\subsection{Intercaler un nombre}
Intercaler un nombre, c'est en trouver un qui se situe \solution{entre} deux autres.
Entre deux nombres décimaux différents, on peut toujours en intercaler une \solution{infinité} !\\

\begin{examplebox}
	Intercaler un nombre entre 4,2 et 4,3.
	\textbf{Astuce :} On ajoute un \solution{zéro} : on cherche un nombre entre 4,20 et 4,30.
	On peut choisir 4,21 ; 4,25 ; 4,29... \\
	Par exemple : \solution[5cm]{$4,2 < 4,25 < 4,3$}.
\end{examplebox}

\section{Placer et repérer sur une demi-droite graduée}

\begin{methodebox}
	\textbf{Méthode pour lire ou placer un point :}
	\begin{enumerate}
		\item \textbf{Observer l'unité} (par exemple de 3 à 4) et compter en combien de parts \solution{égales} elle est partagée.
		\item \textbf{Calculer la valeur d'une graduation.} Si l'unité est partagée en 10, chaque graduation vaut \solution{0,1}. Si elle est partagée en 100, chaque graduation vaut \solution{0,01}.
		\item \textbf{Compter les graduations} depuis un repère connu pour lire ou placer le point.
	\end{enumerate}
\end{methodebox}

\begin{examplebox}
	Sur cette droite, l'unité est partagée en 10. Chaque graduation vaut 0,1.
	\begin{center}
		\begin{tikzpicture}[scale=1.2]
			\draw[thick, ->] (0,0) -- (11,0);
			\foreach \x in {0,10} {
				\draw (\x,-0.2) -- (\x,0.2);
				\node[below] at (\x,-0.4) {\ifnum\x=0 3\else 4\fi};
			}
			\foreach \x in {1,...,9} {\draw (\x,-0.1) -- (\x,0.1);}
			\fill[blue] (7,0) circle (0.08); \node[above] at (7,0.3) {$A$};
			\fill[red] (2,0) circle (0.08); \node[above] at (2,0.3) {$B$};
		\end{tikzpicture}
	\end{center}
	Le point A est à 7 graduations après 3, donc son abscisse est \textbf{\solution{3,7}}. On note $A(3,7)$.
	Le point B a pour abscisse \textbf{\solution{3,2}}. On note $B(3,2)$.
\end{examplebox}

% --- Les exercices et le résumé sont volontairement laissés tels quels ---

\clearpage
\section{Rang d'un chiffre}

Chaque position a un nom et une valeur. Exemple avec $3\,941{,}045$ :
\begin{center}
	\begingroup
	\setlength{\tabcolsep}{3pt} % espace horizontal entre colonnes (local)
	\renewcommand{\arraystretch}{1.15} % hauteur de ligne
	\small % réduit légèrement la taille de police
	\begin{tabularx}{\textwidth}{|*{8}{>{\centering\arraybackslash}X|}}
		\hline
		\textbf{Milliers} & \textbf{Centaines} & \textbf{Dizaines} & \textbf{Unités} &
		\textbf{Dixièmes} & \textbf{Centièmes} & \textbf{Millièmes} & \textbf{Dix-millièmes} \\
		\hline
		\solution{3} & \solution{9} & \solution{4} & \solution{1} &
		\solution{0} & \solution{4} & \solution{5} & \\
		\hline
	\end{tabularx}
	\endgroup
\end{center}


Ainsi, le \textbf{chiffre des dizaines} vaut \solution{4} et le \textbf{chiffre des centièmes} vaut \solution{4}.

\section{Lien avec les unit\'es de mesure}

\begin{definitionbox}
\textbf{Longueurs principales}
\begin{itemize}[label = \textbullet]
    \item $1\,\text{m} = \solution{10}\,\text{dm} = \solution{100}\,\text{cm} = \solution{1\,000}\,\text{mm}$
    \item $1\,\text{km} = \solution{1\,000}\,\text{m}$ ; $1\,\text{hm} = \solution{100}\,\text{m}$ ; $1\,\text{dam} = \solution{10}\,\text{m}$
    \item $1\,\text{dm} = \solution{0{,}1}\,\text{m}$ ; $1\,\text{cm} = \solution{0{,}01}\,\text{m}$ ; $1\,\text{mm} = \solution{0{,}001}\,\text{m}$
\end{itemize}
\end{definitionbox}

\begin{examplebox}
\textbf{Conversions}
\begin{itemize}[label = \textbullet]
    \item $45\,\text{cm} = \solution{0{,}45}\,\text{m}$ ; $3{,}5\,\text{m} = \solution{350}\,\text{cm}$
    \item $1{,}2\,\text{km} = \solution{1\,200}\,\text{m}$ ; $1\,250\,\text{m} = \solution{1{,}25}\,\text{km}$
\end{itemize}
\end{examplebox}

\begin{remarkbox}
De même, on peut faire le lien avec les \textbf{masses} (kg, g, mg) et les \textbf{contenances} (L, cL, mL).
\end{remarkbox}

\clearpage
\section{Valeurs approch\'ees et encadrement}

\begin{definitionbox}
\textbf{Amplitude d'un encadrement} : si $a < x < b$, alors l'amplitude vaut $\solution{b-a}$.
\end{definitionbox}

\begin{examplebox}
\textbf{Exemple} : $10 < 17{,}6 < 20$ est un encadrement d'amplitude \solution{10} de $17{,}6$.
\end{examplebox}

\begin{proprietebox}[Arrondir]
Pour arrondir un nombre d\'ecimal :
\begin{itemize}[label = \textbullet]
    \item au \textbf{dixième} : on regarde le \textbf{centième} ; \solution{$\ge 5$} \,$\Rightarrow$ on augmente le dixième, sinon on laisse.
    \item au \textbf{centième} : on regarde le \textbf{millième} ; \solution{$\ge 5$} \,$\Rightarrow$ on augmente le centième, sinon on laisse.
\end{itemize}
\end{proprietebox}

\begin{examplebox}
\textbf{Valeurs approch\'ees de $3{,}538$} :
\begin{center}
\begin{tabular}{|c|c|}
\hline
\textbf{\`A l'unité près} & \solution{4} \\
\hline
\textbf{Au dixième près} & \solution{3{,}5} \\
\hline
\textbf{Au centième près} & \solution{3{,}54} \\
\hline
\end{tabular}
\end{center}
\end{examplebox}

\section{Exercices d'application}

\begin{exercisebox}
	\textbf{Exercice 1 :} Écrire sous forme de fraction décimale puis sous forme décimale
	\begin{enumerate}
		\item Trois dixièmes : \solution[3cm]{$\frac{3}{10} = 0,3$}
		\item Vingt-sept centièmes : \solution[3cm]{$\frac{27}{100} = 0,27$}
		\item Cent quarante-cinq millièmes : \solution[4cm]{$\frac{145}{1000} = 0,145$}
		\item Douze unités et sept dixièmes : \solution[4cm]{$12 + \frac{7}{10} = 12,7$}
	\end{enumerate}

	\textbf{Exercice 2 :} Décomposer les nombres suivants
	\begin{enumerate}
		\item $34{,}56 = \solution[1cm]{34} + \dfrac{\solution[1cm]{56}}{100}$
		\item $125{,}8 = \solution[1.5cm]{125} + \dfrac{\solution[1cm]{8}}{10}$
		\item $7{,}403 = \solution[1cm]{7} + \dfrac{\solution[1.5cm]{403}}{1~000}$
	\end{enumerate}

	\textbf{Exercice 3 :} Comparer les nombres suivants (utiliser $<$, $>$, ou $=$)
	\setlength{\columnsep}{0.8cm} % espace entre colonnes (local)
	\begin{multicols}{2}
		\begin{enumerate}[itemsep=0.2em, topsep=0.2em, leftmargin=*, label=\arabic*)]
			\item $12{,}3 \solution[8mm]{=} 12{,}30$
			\item $5{,}67 \solution[8mm]{<} 5{,}7$
			\item $8{,}09 \solution[8mm]{<} 8{,}1$
			\item $15{,}999 \solution[8mm]{<} 16$
		\end{enumerate}
	\end{multicols}

	\textbf{Exercice 4 :} Ranger dans l'ordre croissant les nombres : $2{,}1$ ; $2{,}01$ ; $2{,}11$ ; $2{,}101$.
	
	\solution[6cm]{$2,01 < 2,1 < 2,101 < 2,11$}

	\textbf{Exercice 5 :} Encadrer 7,384
	\begin{enumerate}
		\item à l'unité : \solution[1cm]{7} $< 7{,}384 <$ \solution[1cm]{8}
		\item au dixième : \solution[1.5cm]{7,3} $< 7{,}384 <$ \solution[1.5cm]{7,4}
		\item au centième : \solution[1.5cm]{7,38} $< 7{,}384 <$ \solution[1.5cm]{7,39}
	\end{enumerate}
\end{exercisebox}

\vspace{0.75cm}

\begin{center}
	\fcolorbox{black}{lightgray!20}{%
		\begin{minipage}{0.9\textwidth}
			\textbf{À retenir absolument :}
			\begin{itemize}[label = $\bullet$]
				\item Une \textbf{fraction décimale} a pour dénominateur \solution{10, 100, 1~000...}
				\item Un \textbf{nombre décimal} peut s'écrire avec une \solution{virgule} ou comme une \solution{fraction décimale}.
				\item On peut \textbf{ajouter ou enlever des zéros} à la fin de la partie décimale sans changer la valeur du nombre ($3,5 = 3,50$).
				\item Pour \textbf{comparer}, on regarde d'abord la partie \solution{entière}, puis les \solution{dixièmes}, les \solution{centièmes}...
				\item Sur une \textbf{droite graduée}, il faut toujours trouver la valeur d'une seule \solution{graduation}.
			\end{itemize}
		\end{minipage}
	}
\end{center}
